\section{Research Methodology}
\label{sec:methodology}

\subsection{Study design}
\label{subsec:study-design}

This study evaluates a cross-shard commit protocol through bounded property
verification, mutation-based property validation, bounded implementation-model
trace conformance, and verification-cost characterization. The four methods
address complementary questions rather than serving as interchangeable
evidence.

The design was frozen before the definitive experimental campaign. Research
questions, hypotheses, factors, configuration profiles, seeds, repetitions,
resource limits, statistical procedures, and rules for incomplete executions
were specified before observing the final results. Changes to these elements
after inspection of the definitive outcomes would require a new protocol
version.

\subsection{Research questions and hypotheses}
\label{subsec:research-questions}

\paragraph{RQ1: Preservation of declared properties.}

Does the cross-shard commit protocol preserve the declared properties under the
explored interleavings within the documented bounds?

The corresponding hypothesis is:

\begin{quote}
\textbf{H1.} Valid models produce no property violations within the declared
bounds.
\end{quote}

RQ1 is interpreted only for the executed property, tool, configuration, and
bound, not as an unbounded correctness proof.

\paragraph{RQ2: Defect-detection capability.}

Do TLC and Alloy detect predefined scientific mutations that remove controls
from the cross-shard protocol?

The corresponding hypothesis is:

\begin{quote}
\textbf{H2.} Each scientific mutant produces a violation of its designated
target property.
\end{quote}

Detection requires the designated target property to be reported among the
violations, although a mutant may violate additional properties.

\paragraph{RQ3: Bounded implementation-model trace conformance.}

Do observable traces produced by the Java implementation correspond to action
sequences admitted by the TLA+ model within the evaluated scenarios and bounds?

The corresponding hypothesis is:

\begin{quote}
\textbf{H3.} Valid traces are accepted and corrupted traces are rejected at the
expected point.
\end{quote}

RQ3 provides bounded implementation-model trace-conformance evidence rather
than a general refinement proof or behavioral equivalence.

\paragraph{RQ4: Verification cost.}

How does verification cost change as the number of shards, concurrent
transfers, bounds, and enabled faults increase?

The preregistered hypothesis is:

\begin{quote}
\textbf{H4.} Model-checking cost grows nonlinearly with the number of concurrent
transfers and shards.
\end{quote}

H4 is empirical rather than assumed and is evaluated from observed
profile-level cost measurements.

\subsection{Methodological layers}
\label{subsec:methodological-layers}

The four questions form three verification and validation layers plus one
cost-analysis layer.

\subsubsection{Bounded property-verification layer}

RQ1 evaluates the valid TLA+ and Alloy specifications against the seven
protocol properties in Section~\ref{subsec:properties}. The experimental unit
consists of:

\begin{itemize}
\item a protocol property
\item a formal tool
\item a bounded configuration
\item a measured repetition
\end{itemize}

Response variables include property outcome, available state-space information,
exploration depth or scope, elapsed time, maximum resident memory, exit status,
and counterexample information. A completed execution without a counterexample
provides evidence only within its declared finite configuration. Timeout and
out-of-memory outcomes remain incomplete or censored observations.

\subsubsection{Mutation-based property validation layer}

RQ2 evaluates scientific mutants that remove specific protocol controls:
replay protection, receipt-dependent destination credit, terminal-decision
consistency, timeout release of source funds, and quorum enforcement. The
experimental unit consists of:

\begin{itemize}
\item a scientific mutant
\item a formal tool
\item a bounded configuration
\item a measured repetition
\end{itemize}

Each mutant has one predefined target property, and detection is credited only
when that property is violated. The aggregate measure is

\begin{equation}
\mathrm{MutationScore}
=
\frac{N_{\mathrm{detected}}}
{N_{\mathrm{mutants}}},
\label{eq:mutation-score}
\end{equation}

where \(N_{\mathrm{detected}}\) is the number of predefined mutants detected
through their target properties and \(N_{\mathrm{mutants}}\) is the number
evaluated. The score applies only to the predefined catalogue and does not
establish property-suite completeness.

\subsubsection{Bounded implementation-model trace-conformance layer}

RQ3 evaluates deterministic Java traces projected through the Java-to-TLA+
abstraction in Section~\ref{subsec:formal-abstraction}. The experimental unit
consists of:

\begin{itemize}
\item a valid scenario or negative mutation
\item a deterministic seed
\item the corresponding TLC replay outcome
\end{itemize}

Two trace classes are evaluated: valid traces generated from supported protocol
scenarios and corrupted traces containing predefined structural or semantic
defects. A valid trace is classified successfully when its projected action
sequence is admitted by the replay model. A corrupted trace succeeds
experimentally when it is rejected at the predefined diagnostic point.

For negative traces, the observed diagnostic must also agree with the expected
abstract step, concrete step, action, and transfer identifier. These criteria
distinguish generic replay failure from rejection at the intended protocol
location. The evidence remains bounded to the frozen scenario, mutation, seed,
and projection catalogue.

\subsubsection{Verification-cost characterization layer}

RQ4 characterizes bounded verification cost. The experimental unit consists of:

\begin{itemize}
\item a formal tool
\item a bound profile
\item a fault profile when applicable
\item a measured repetition
\end{itemize}

Response variables are elapsed time, maximum resident memory, generated and
distinct states when available, exploration depth or Alloy scope, timeout,
out-of-memory outcome, and within-tool growth ratios. Because TLC and Alloy
expose different exploration semantics and state representations, cost trends
are analyzed within each tool rather than by treating their absolute
measurements as interchangeable.

\subsection{Triangulation of evidence}
\label{subsec:evidence-triangulation}

The layers address different weaknesses in the scientific argument. RQ1
evaluates bounded property outcomes but cannot show by itself that the
properties expose relevant defects. RQ2 adds targeted defect sensitivity.
Formal-model evidence alone does not establish agreement with the executable
implementation, so RQ3 adds bounded trace conformance with positive and
negative cases. RQ4 characterizes the computational envelope in which the
formal evidence is obtained.

The dependency among the layers is deliberate. Successful bounded checking in
RQ1 is more informative when RQ2 shows that selected weakened controls trigger
their target properties. Those two layers remain statements about formal
representations, so RQ3 confronts the abstraction with executable traces. RQ4
then reports the resource envelope within which the formal evidence was
obtained.

Thus, the first three layers provide complementary evidence about bounded
property outcomes, defect sensitivity, and implementation-model trace
conformance, while the fourth characterizes their verification cost. The layers
strengthen different parts of the argument, but no layer is interpreted as
replacing the limitations of another.

\subsection{Treatment of incomplete executions}
\label{subsec:incomplete-executions}

Experimental outcomes are retained according to observed status.
Counterexamples, timeouts, out-of-memory results, and other scientifically
meaningful unfavorable outcomes are not removed. A failed instrumentation
attempt may be repeated only under the predefined protocol rules, not because a
run produced an unexpected scientific result.

Timeout and out-of-memory observations are retained as censored results rather
than assigned artificial elapsed-time or memory values or reclassified as
successful verification. This prevents incomplete runs from being converted
into positive evidence.

\subsection{Interpretation boundaries}
\label{subsec:method-boundaries}

All claims remain bounded by the declared protocol model, properties,
configurations, scopes, scenarios, seeds, mutant catalogue, and execution
environment. Absence of a counterexample is not an unrestricted proof,
mutation analysis does not establish property-suite completeness, and trace
replay does not establish general refinement or behavioral equivalence. TLC and
Alloy have different semantics and cost units, so their absolute measurements
do not establish intrinsic tool superiority. The evaluated configurations are
abstractions rather than unrestricted production blockchain systems.
